
Our work extends code generation benchmarks with process metrics measuring debugging efficiency.

\subsubsection{Code Generation Benchmarks}

HumanEval and MBPP established pass@k as the standard correctness metric for code generation, measuring the probability that at least one of k samples passes unit tests. These benchmarks evaluate generation quality but not iterative debugging---problems provide 1-5 test cases, insufficient to reveal strategic behavior. HumanEval+ improves test coverage to detect brittleness but retains outcome-focused evaluation without analyzing debugging trajectories.

Competitive programming benchmarks (CodeContests, APPS) use multi-test problems (10-100 test cases) where strategic debugging becomes measurable. Existing work reports final solve rates or test-pass percentages, not how efficiently agents reach these outcomes. Our framework analyzes debugging trajectories on multi-test problems, measuring fix-impact-ratio and error clustering---metrics invisible to total pass rate.

\subsubsection{Execution Feedback in Code Generation}

Execution-based feedback improves code generation through reinforcement learning or iterative refinement with error messages. These approaches use feedback for training (e.g., RL reward signals from test execution) rather than evaluation---they measure whether feedback improves generation, not how agents strategically utilize feedback during debugging.

Our work is orthogonal: we measure strategic debugging behavior (clustering errors, prioritizing fixes) as an evaluation framework, applicable to any agent regardless of training method. An RL-trained agent and a prompted GPT-4 agent can both be evaluated with fix-impact-ratio to quantify debugging efficiency---the metric measures process, not training paradigm.

\subsubsection{Multi-Test Programming Evaluation}

Multi-test problems provide 15-50 test cases per problem, but existing work focuses on final correctness rates. SWE-bench evaluates GitHub issue resolution with repository-level context, measuring patch success without analyzing debugging efficiency when test suites are available.

Our framework could extend to SWE-bench scenarios where test suites exist---measuring how efficiently agents pass tests reveals strategic capability beyond issue resolution rates. Similarly, multi-test competitive programming benchmarks support fix-impact-ratio measurement, but prior work hasn't analyzed debugging trajectories systematically.

\subsubsection{Positioning}

We extend outcome-focused benchmarks (HumanEval pass@k, solve rates) by adding process metrics revealing how agents debug. Our metrics are execution-based (no subjective model judges required), applicable to any multi-test programming scenario, and orthogonal to training methods (measure behavior, not training approach). This positions our work as a complementary evaluation dimension: correctness (pass@k) + efficiency (fix-impact-ratio) together provide a fuller picture of agent capability.

The key distinction is that we measure debugging as a process with testable stages (clustering $\rightarrow$ prioritization), validated via execution-based metrics rather than relying on final pass rates alone. This reveals strategic debugging operates as a feedback loop effective within revealed test execution, but not a learning system enabling pattern transfer to unseen tests---a boundary existing benchmarks cannot detect.
